\section{仿真结果}
值得注意的是，所提方法在各种场景下均显著提升了跟踪精度。然而，传统方法难以有效解决这一问题。因此，本文方法具有重要意义。此外，实验充分验证了方法的有效性。总之，该方法为后续研究奠定了坚实基础。

图~\ref{fig:a} 给出了各方法的对比结果。所提方法优于基线方法。该方法有效地改善了低信噪比条件下幅度证据不足导致的关联歧义增加问题。
